\documentclass[10pt,a4paper]{amsart}
\usepackage{amsmath,amssymb,amsthm,graphicx}
\usepackage{verbatim}
\newenvironment{reference}{\comment}{\endcomment}
%\newenvironment{reference}{}{}
\newenvironment{slogan}{\comment}{\endcomment}
\newenvironment{history}{\comment}{\endcomment}

% For commutative diagrams we use Xy-pic
\usepackage[all]{xy}

% We use 2cell for 2-commutative diagrams.
\xyoption{2cell}
% Theorem environments.
%
\theoremstyle{plain}
\newtheorem{theorem}[subsection]{Theorem}
\newtheorem{proposition}[subsection]{Proposition}
\newtheorem{lemma}[subsection]{Lemma}

\theoremstyle{definition}
\newtheorem{definition}[subsection]{Definition}
\newtheorem{example}[subsection]{Example}
\newtheorem{exercise}[subsection]{Exercise}
\newtheorem{situation}[subsection]{Situation}

\theoremstyle{remark}
\newtheorem{remark}[subsection]{Remark}
\newtheorem{remarks}[subsection]{Remarks}

\numberwithin{equation}{subsection}

% Macros
%
\def\lim{\mathop{\mathrm{lim}}\nolimits}
\def\colim{\mathop{\mathrm{colim}}\nolimits}
\def\Spec{\mathop{\mathrm{Spec}}}
\def\Hom{\mathop{\mathrm{Hom}}\nolimits}
\def\Ext{\mathop{\mathrm{Ext}}\nolimits}
\def\SheafHom{\mathop{\mathcal{H}\!\mathit{om}}\nolimits}
\def\SheafExt{\mathop{\mathcal{E}\!\mathit{xt}}\nolimits}
\def\Sch{\mathit{Sch}}
\def\Mor{\mathop{\mathrm{Mor}}\nolimits}
\def\Ob{\mathop{\mathrm{Ob}}\nolimits}
\def\Sh{\mathop{\mathit{Sh}}\nolimits}
\def\NL{\mathop{N\!L}\nolimits}
\def\CH{\mathop{\mathrm{CH}}\nolimits}
\def\proetale{{pro\text{-}\acute{e}tale}}
\def\etale{{\acute{e}tale}}
\def\QCoh{\mathit{QCoh}}
\def\Ker{\mathop{\mathrm{Ker}}}
\def\Im{\mathop{\mathrm{Im}}}
\def\Coker{\mathop{\mathrm{Coker}}}
\def\Coim{\mathop{\mathrm{Coim}}}

% Boxtimes
%
\DeclareMathSymbol{\boxtimes}{\mathbin}{AMSa}{"02}

%
% Macros for moduli stacks/spaces
%
\def\QCohstack{\mathcal{QC}\!\mathit{oh}}
\def\Cohstack{\mathcal{C}\!\mathit{oh}}
\def\Spacesstack{\mathcal{S}\!\mathit{paces}}
\def\Quotfunctor{\mathrm{Quot}}
\def\Hilbfunctor{\mathrm{Hilb}}
\def\Curvesstack{\mathcal{C}\!\mathit{urves}}
\def\Polarizedstack{\mathcal{P}\!\mathit{olarized}}
\def\Complexesstack{\mathcal{C}\!\mathit{omplexes}}
% \Pic is the operator that assigns to X its picard group, usage \Pic(X)
% \Picardstack_{X/B} denotes the Picard stack of X over B
% \Picardfunctor_{X/B} denotes the Picard functor of X over B
\def\Pic{\mathop{\mathrm{Pic}}\nolimits}
\def\Picardstack{\mathcal{P}\!\mathit{ic}}
\def\Picardfunctor{\mathrm{Pic}}
\def\Deformationcategory{\mathcal{D}\!\mathit{ef}}
\usepackage{cleveref}
\crefname{originalreference}{original reference}{original references}
\Crefname{originalreference}{Original reference}{Original references}
\makeatletter
\newlabel{lemma-surjective-cotangent-space}{{06GZ}{source}{}{sourcepointer-1}{}}
\newlabel{lemma-surjective-cotangent-space@cref}{{[originalreference][0][]06GZ}{[1][0][]source}{}{}{}}
\makeatother
\title{Injectivity of the restriction of derivations forces surjectivity of the ring map}
\author{The Stacks Project Authors}
\date{stacks-project:formal-defos.tex:local-snapshot-20261012; GNU FDL 1.2 or later}
\begin{document}
\maketitle
\noindent\textit{Editorial note.} All delivered mathematical text is the author's own, copied from the pinned Stacks Project chapter file formal-defos.tex; no statement or proof step was rewritten, reordered or removed.\par\medskip
\noindent\textit{Editorial note.} The counted proof is the author's complete argument: the translation of nonsurjectivity into the nonvanishing of the relative cotangent quotient, the explicit verification of the Leibniz rule for the quotient map, and the deduction of a nonzero derivation trivial on the source ring.\par\medskip
\noindent\textit{Editorial note.} The label lemma-derivations-surjective names the counted statement, and the proof environment immediately following it in the source is delivered verbatim as its proof.\par\medskip
\noindent\textit{Editorial note.} The counted statement is carried by the author's own lemma environment in the chapter source.\par\medskip
\noindent\textit{Editorial note.} Pre-open gate: python operations/agent-authoring/tools/precard\_gate.py --source formal-defos.tex --label lemma-derivations-surjective --start 1219 --end 1248 --local-id 28025 returned RESULT=OK (L28025 free, occupancy FREE, public interval clear). Candidate taken from sources/stacks-formal-defos/candidates-halfA.json.\par\medskip
\noindent\textit{Editorial note.} The statement and proof are copied from the pinned chapter file; the cited inputs are kept as external pointers with their canonical Stacks tags.\par\medskip
\noindent\textit{Allowed standard prerequisites (editorial).}\begin{itemize}
\item Detecting surjectivity on the relative cotangent space. Same chapter, the lemma that a map in the completed category is surjective exactly when its relative cotangent space vanishes. Supplies the nonvanishing of the relative cotangent space when f is not surjective, hence the nonvanishing of the quotient Q.
\item Structure of the quotient by the subring and the square of the maximal ideal. Standard properties of the quotient of a local algebra by the image of a ring map plus the square of its maximal ideal, making it a vector space over the residue field. Supplies the identification of Q as a k-vector space and as an S-module through the residue field.
\end{itemize}\medskip
\ifdefined\hypertarget\hypertarget{sourcepointer-1}{}\fi
\noindent\textit{Source reference.} Original reference 06GZ: sources/stacks-formal-defos/pinned-original-source.tex, source lines 1086--1086. Prerequisite pointer invoked at proof line(s) 1229; it names a standard fact of the same chapter that the counted argument uses, and every step of the counted proof built on it is retained verbatim in the delivered body.\par\medskip
\begin{lemma}
\label{lemma-derivations-surjective}
Let $f : R \to S$ be a morphism of $\widehat{\mathcal{C}}_\Lambda$.
If $\text{Der}_\Lambda(S, k) \to \text{Der}_\Lambda(R, k)$ is injective,
then $f$ is surjective.
\end{lemma}

\begin{proof}
If $f$ is not surjective, then
$\mathfrak m_S/(\mathfrak m_R S + \mathfrak m_S^2)$ is nonzero by
Lemma \ref{lemma-surjective-cotangent-space}.
Then also $Q = S/(f(R) + \mathfrak m_R S + \mathfrak m_S^2)$ is nonzero.
Note that $Q$ is a $k = R/\mathfrak m_R$-vector space via $f$. We turn
$Q$ into an $S$-module via $S \to k$. The quotient
map $D : S \to Q$ is an $R$-derivation: if $a_1, a_2 \in S$, we can write
$a_1 = f(b_1) + a_1'$ and $a_2 = f(b_2) + a_2'$ for some
$b_1, b_2 \in R$ and $a_1', a_2' \in \mathfrak m_S$. Then
$b_i$ and $a_i$ have the same image in $k$ for $i = 1, 2$ and
\begin{align*}
a_1a_2 & = (f(b_1) + a_1')(f(b_2) + a_2') \\
& = f(b_1)a_2' + f(b_2)a_1' \\
& = f(b_1)(f(b_2) + a_2') + f(b_2)(f(b_1) + a_1') \\
& = f(b_1)a_2 + f(b_2)a_1
\end{align*}
in $Q$ which proves the Leibniz rule. Hence $D : S \to Q$ is a
$\Lambda$-derivation which is zero on composing with $R \to S$.
Since $Q \not = 0$ there also exist derivations $D : S \to k$ which
are zero on composing with $R \to S$, i.e.,
$\text{Der}_\Lambda(S, k) \to \text{Der}_\Lambda(R, k)$ is not injective.
\end{proof}

\end{document}
